\section{A hierarchy of branching lattice points}
\label{sec:hierarchy}

The preceding block estimates produce a function whose mean is close to one,
although the function may vanish at many targets.  We now arrange blocks in
decreasing sizes.  The bits in one block sample the mean of the preceding
block; a final block of bounded size supplies a starting point.  Keeping all
the intermediate branches is what preserves their total weight.
The data are chosen from the largest block to the smallest: each new
block samples the function already fixed in its predecessor. For a given
target, the lattice points are selected in the reverse order, starting
with the terminal block and then choosing the earlier blocks conditional
on the suffix already selected.

\subsection{The block sizes and the lattice}

Fix a sufficiently large absolute integer $B$, whose requirements will be
specified below.  Starting with an integer $b_1\ge B$, form a finite list by
putting
\[
 b_{i+1}=\lceil b_i^{0.9}\rceil
\]
whenever this number is at least $B$, and otherwise stopping.  Let $r$ be the
last index.  Increasing $B$ ensures that
\begin{equation}
 b_{i+1}\le \tfrac12 b_i,\qquad
 B\le b_r<B^{1/0.9},\qquad
 \sum_{i=1}^r b_i\le 2b_1.
 \label{eq:hierarchy-sizes}
\end{equation}
Use the notation of the preceding section at size $b_i$:
$h_i=h_{b_i}$, $p_i=p_{b_i}$, $G_i=\F_{p_i}^{b_i}$, $F_i=F_{b_i}$, and
$P_{i,t}=P_{b_i,t}$.  The eligible and usable anchors are always represented
by $t_z=z/p_i\in[0,1)^{b_i}$.

Apply Lemma~\ref{lem:usable} successively.  In the first block there is no
preceding function.  In each subsequent block apply it to $g_{i-1}$ on
$G_{i-1}$.  Its hypotheses hold for sufficiently large $B$: if $b=b_i$,
then $b_{i-1}\le b^{1/0.9}$ and
\[
 5b_{i-1}^{0.56}\le b^{0.65},\qquad
 \log|G_{i-1}|=b_{i-1}\log p_{i-1}
       =O(b_{i-1}^2\log b_{i-1})\le b^4.
\]
The first inequality uses $0.56/0.9<0.65$, and the second has the slack
$2/0.9<4$.  The means belong to $[1/2,2]$ after another increase of $B$.
We obtain usable sets $U_i\subset G_i$, lines $C_i\subset G_i$, and, for
$i<r$, matrices $W_i$ of size $b_i$ by $b_{i+1}$ with entries in $\F_{p_i}$.
Thus
\begin{equation}
 \begin{split}
 g_i(x)&=\max_{c\in C_i}
       \frac{F_i(t_{x-c})}{p_i}\ind_{U_i}(x-c),\\
 |\bar g_i-1|&\le A b_i^{-1/100},\qquad
 0\le g_i\le \exp(5b_i^{0.56}),
 \end{split}
 \label{eq:hierarchy-functions}
\end{equation}
where $\bar g_i=\E_{x\in G_i}g_i(x)$ and $A$ is absolute.  Moreover, every
$z\in U_i$ with $i>1$ satisfies
\begin{equation}
 \sum_{\epsilon\in\{0,1\}^{b_i}} P_{i,t_z}(\epsilon)
       g_{i-1}(x-W_{i-1}\epsilon)
 \ge (1-A b_i^{-1/100})\bar g_{i-1}
 \quad (x\in G_{i-1}).
 \label{eq:hierarchy-biased-sampler}
\end{equation}
All constants in these estimates are independent of the number of blocks.

For $i\le r$, define the raw coordinate lattice
\begin{equation}
 L_i=\Z^{b_i}+p_i^{-1}C_i.
 \label{eq:hierarchy-raw-lattice}
\end{equation}
Here the right side means the union of the integer translates of the
standard representatives divided by $p_i$; it is independent of the choice
of representatives.  It is a lattice of determinant $p_i^{-1}$.  Indeed,
choose a generator $c$ of $C_i$ with its $j$th coordinate equal to $1$.
The vector $c/p_i$ and the vectors $e_k$ for $k\ne j$ are a basis:
they generate $e_j=p_i(c/p_i)-\sum_{k\ne j}c_k e_k$ and all of $L_i$,
and their determinant has absolute value $1/p_i$.

Adjoin a terminal block of size
\[
 T=b_{r+1}=b_r\lceil\log_2 p_r\rceil,
 \qquad h_{r+1}=1,\qquad L_{r+1}=\Z^T.
\]
Its size is bounded by a constant depending only on the fixed cutoff $B$.
Define $W_r$, of size $b_r$ by $T$, by using the columns
$2^k e_j\pmod {p_r}$ for each $1\le j\le b_r$ and
$0\le k<\lceil\log_2 p_r\rceil$.  Consequently
\begin{equation}
 \{W_r\epsilon:\epsilon\in\{0,1\}^T\}=G_r.
 \label{eq:hierarchy-terminal-exhaustion}
\end{equation}
Use standard integer representatives for every $W_i$, and put
\[
 A_i=p_i^{-1}W_i\quad (1\le i\le r),\qquad
 L=\prod_{i=1}^{r+1}L_i,\qquad D=\sum_{i=1}^{r+1}b_i.
\]
Each entry of $A_i$ belongs to $[0,1)$ and is an integer multiple of
$1/p_i$.  On raw block coordinates $\ell=(\ell_1,\ldots,\ell_{r+1})$,
define the invertible block upper triangular map $M$ by
\begin{equation}
 (M\ell)_i=h_i(\ell_i+A_i\ell_{i+1})\quad (i\le r),
 \qquad (M\ell)_{r+1}=\ell_{r+1}.
 \label{eq:hierarchy-map}
\end{equation}
In particular, $\Lambda_v=ML$ is a single full-rank lattice, with
\begin{equation}
 D_v=\det\Lambda_v=\prod_{i=1}^r\frac{h_i^{b_i}}{p_i}.
 \label{eq:hierarchy-determinant}
\end{equation}
The fractional off-diagonal entries do not change this conclusion: $M$ is
one invertible linear map applied to the product lattice $L$.

\begin{proposition}[Weighted vertical patterns]
\label{prop:vertical}
There are absolute constants $B,C,c,\kappa>0$ such that, for every integer
$b_1\ge B$, the construction above has
\begin{equation}
 D\le C(b_1+1),\qquad
 R^2:=\sum_{i=1}^{r+1}b_i h_i^2
       \le C\bigl(b_1\log(b_1+1)+1\bigr).
 \label{eq:hierarchy-dimension-radius}
\end{equation}
For every $y\in\R^D$ there is a finite nonempty pattern
$\mathcal P_y\subset L$ such that
\begin{align}
 \|(y-M\ell)_i\|_\infty&\le h_i
       \qquad(\ell\in\mathcal P_y,\ 1\le i\le r+1),
       \label{eq:hierarchy-residual}\\
 \sum_{\ell\in\mathcal P_y}D_v\gamma_D(y-M\ell)&\ge\kappa,
       \label{eq:hierarchy-total}\\
 \max_{\ell\in\mathcal P_y}D_v\gamma_D(y-M\ell)
       &\le \exp(-c b_1^{0.70}).
       \label{eq:hierarchy-cap}
\end{align}
The terminal block is constant on the pattern.  For each ordinary block
$i$ and each fixed suffix $(\ell_{i+1},\ldots,\ell_{r+1})$, its retained
values belong to $\lambda_i+\{0,1\}^{b_i}$ for one
$\lambda_i\in L_i$ depending on that suffix.  In particular
$|\mathcal P_y|\le 2^D$.

Order the scalar coordinates by increasing block index and denote them by
$\ell[1],\ldots,\ell[D]$.  Conditional on any fixed scalar suffix
$\ell[j+1],\ldots,\ell[D]$, the possible next values $\ell[j]$ are at most
two; if there are two, they differ by $1$.  This property is preserved upon
deleting any subset of the pattern.

If $y\in M[0,1)^D$, all raw coordinates can in addition be restricted to
the fixed finite alphabets
\begin{equation}
 \begin{aligned}
 \ell_i&\in\bigl(p_i^{-1}\Z\cap[-Q_{\rm raw},Q_{\rm raw}]\bigr)^{b_i}
         &&(i\le r),\\
 \ell_{r+1}&\in\bigl(\Z\cap[-Q_{\rm raw},Q_{\rm raw}]\bigr)^T,
         &&Q_{\rm raw}=1+(D+2)^D.
 \end{aligned}
 \label{eq:hierarchy-alphabets}
\end{equation}
For an ordinary block the logarithm of each scalar alphabet size is
$O(D\log(D+2)+b_i\log b_i)$; the terminal version omits the second term.
All implicit constants are absolute.
\end{proposition}

\begin{proof}
The size and radius bounds follow from Equation~\eqref{eq:hierarchy-sizes},
the formula $h_i^2=(9/4)\log b_i$, and the absolute bound on $T$.  We prove
the weight estimates by an induction that allows an empty intermediate
pattern.

\paragraph{Truncated patterns.}
Retain just the ordinary blocks $1,\ldots,i$, omitting the interaction of
block $i$ with its successor.  Write $M^{[i]}$ for this truncated map,
$d_i=\sum_{j\le i}b_j$, and
$D_v^{[i]}=\prod_{j\le i}h_j^{b_j}/p_j$.  We claim that for every
$y=(y_1,\ldots,y_i)\in\R^{d_i}$ there is a pattern with total weight at
least
\begin{equation}
 \sum_{\ell\in\mathcal P^{[i]}_y}
      D_v^{[i]}\gamma_{d_i}(y-M^{[i]}\ell)
 \ge a_i g_i(x_i),\qquad
 x_i=\lfloor p_i y_i/h_i\rfloor\pmod {p_i},
 \label{eq:hierarchy-induction}
\end{equation}
where
\begin{equation}
 a_i\ge\prod_{j=1}^i(1-C_0 b_j^{-1/100})
 \label{eq:hierarchy-product}
\end{equation}
for an absolute $C_0$.  The residual and block-branching properties in the
proposition are part of this induction.  Empty patterns are permitted when
$g_i(x_i)=0$.

Fix the lexicographic order on the standard representatives of each
finite group; when a maximum has ties, choose the least maximizing anchor.
Suppose $g_i(x_i)>0$.  Choose a usable anchor $z\in x_i-C_i$ realizing
the maximum in Equation~\eqref{eq:hierarchy-functions}.  By
Lemma~\ref{lem:block-cell} there is $\lambda_i\in L_i$ such that
\[
 t=y_i/h_i-\lambda_i\in t_z+[0,1/p_i)^{b_i}.
\]
Consider all choices $\ell_i=\lambda_i+\epsilon$, with
$\epsilon\in\{0,1\}^{b_i}$.  Their individual block weights satisfy
\begin{equation}
 \frac{h_i^{b_i}}{p_i}\gamma_{b_i}(h_i(t-\epsilon))
 \ge e^{-\delta_i}\frac{F_i(t_z)}{p_i}P_{i,t_z}(\epsilon),
 \qquad \delta_i=\frac{b_i h_i^2}{p_i},
 \label{eq:hierarchy-rounded-weight}
\end{equation}
by Equation~\eqref{eq:block-rounding}.  Since each coordinate of $t$ is in
$[0,1)$, the corresponding residual has absolute value at most $h_i$.
For $i=1$, summing Equation~\eqref{eq:hierarchy-rounded-weight} over all
bits proves Equation~\eqref{eq:hierarchy-induction} with $a_1=e^{-\delta_1}$.

For $i>1$, attach the preceding induction pattern to each bit choice,
using the modified target whose last block is
\[
 \widetilde y_{i-1}
    =y_{i-1}-h_{i-1}A_{i-1}(\lambda_i+\epsilon),
\]
and whose earlier blocks are unchanged.  This is allowed even when its
preceding pattern is empty.  The new preceding grid index is exactly
$x-W_{i-1}\epsilon$, where
\begin{equation}
 x=\left\lfloor p_{i-1}y_{i-1}/h_{i-1}
                    -W_{i-1}\lambda_i\right\rfloor\pmod {p_{i-1}}.
 \label{eq:hierarchy-grid-shift}
\end{equation}
Indeed, the vector $W_{i-1}\epsilon$ is integral and hence commutes with
the coordinate floor as a translation.  The fixed vector
$W_{i-1}\lambda_i$ need not be integral; in particular, no compatibility
between the adjacent primes is required.

Patterns attached to different bit vectors are disjoint because their
last blocks differ.  Since $F_i(t_z)/p_i=g_i(x_i)$, multiplying their
weights by Equation~\eqref{eq:hierarchy-rounded-weight} and summing over
every bit vector gives total weight at least
\[
 a_{i-1}e^{-\delta_i}g_i(x_i)
 \sum_{\epsilon\in\{0,1\}^{b_i}}
 P_{i,t_z}(\epsilon)g_{i-1}(x-W_{i-1}\epsilon).
\]
Equation~\eqref{eq:hierarchy-biased-sampler} recovers the preceding mean
from this sum.  Thus we may take
\[
 \begin{aligned}
 a_i&=a_{i-1}e^{-\delta_i}(1-A b_i^{-1/100})\bar g_{i-1}\\
    &\ge a_{i-1}e^{-\delta_i}
       (1-A b_i^{-1/100})(1-A b_{i-1}^{-1/100}),
 \end{aligned}
\]
where the second line uses the mean bound in
Equation~\eqref{eq:hierarchy-functions}.  The prime estimates imply that
$\delta_i$ decays faster than any fixed negative power of $b_i$.  As
$b_{i-1}\ge b_i$, increasing one absolute constant $C_0$ bounds the new
multiplier below by $1-C_0b_i^{-1/100}$ and yields
Equation~\eqref{eq:hierarchy-product}, including its base case.  This step
uses the weighted sum over all bits, not a choice of one successful bit.

\paragraph{Summable losses.}
Set $a=1/100$.  Equation~\eqref{eq:hierarchy-sizes} gives the explicit bound
\begin{equation}
 \sum_{j=1}^r b_j^{-a}
 \le b_r^{-a}\sum_{k=0}^{r-1}2^{-ak}
 \le \frac{B^{-a}}{1-2^{-a}}.
 \label{eq:hierarchy-summability}
\end{equation}
Choose $B$ so that $C_0B^{-a}\le1/2$.  Using
$\log(1-u)\ge-2u$ for $0\le u\le1/2$, the products in
Equation~\eqref{eq:hierarchy-product} are bounded below by the positive
absolute constant
\[
 a_*:=\exp\left(-\frac{2C_0B^{-a}}{1-2^{-a}}\right).
\]
This lower bound is independent of $r$ and $b_1$.

\paragraph{The terminal correction.}
For the full target $y$, choose a terminal raw point of the form
$\ell_{r+1}=\lfloor y_{r+1}\rfloor+\epsilon$.
After accounting for its integer part, varying $\epsilon$ shifts the last
ordinary grid index by $-W_r\epsilon$.  By
Equation~\eqref{eq:hierarchy-terminal-exhaustion} we can place that index at
any point of $G_r$.  Since $\bar g_r\ge1/2$, choose it at a point with
$g_r\ge1/2$.  Apply Equation~\eqref{eq:hierarchy-induction} to the first $r$
blocks with this fixed terminal contribution subtracted from $y_r$.
Every coordinate of the terminal residual has absolute value at most one,
and therefore its Gaussian factor is at least $(2\pi e)^{-T/2}$.
If $T_*(B)$ is an upper bound for all possible terminal dimensions, the
resulting full pattern has total weight at least
\[
 \kappa:=\tfrac12 a_*(2\pi e)^{-T_*(B)/2}>0.
\]
The cutoff $B$ is now fixed, so $\kappa$ is absolute.  This proves
Equation~\eqref{eq:hierarchy-total} and, in particular, nonemptiness.

\paragraph{Individual weights and branching.}
Every retained point comes from a usable anchor in every ordinary block.
The atom estimate in Equation~\eqref{eq:block-atom}, followed by the upper
rounding estimate in Equation~\eqref{eq:block-rounding}, bounds its
individual normalized block weight by
\[
 e^{\delta_i}\exp(-c_0 b_i^{0.70})
 \le \exp(-c b_i^{0.70}),
\]
after increasing $B$ and decreasing the absolute constant $c>0$.
The terminal Gaussian factor is at most $(2\pi)^{-T/2}\le1$.
Multiplication over all blocks proves Equation~\eqref{eq:hierarchy-cap}.
The preceding construction already proves all residual bounds.

Different bit vectors at a recursive step have different last blocks.
Thus a fixed suffix of later blocks identifies a unique recursive branch.
On that branch the construction chose one base $\lambda_i$ before
enumerating its bits and retained a subset of its Boolean translate.
Distinct branches are distinct raw lattice points. Once a scalar suffix is fixed,
the next coordinate can therefore only be the corresponding coordinate
of $\lambda_i$ or that coordinate plus one.  Further deletion of points
cannot destroy this assertion.  This proves the stated scalar suffix
rule and the cardinality bound.

\paragraph{Finite alphabets.}
Finally suppose $y=Mz$ with $z\in[0,1)^D$, and put $e_i=z_i-\ell_i$.
Equation~\eqref{eq:hierarchy-residual} and the entry bounds on $A_i$ imply
\[
 \|e_{r+1}\|_\infty\le1,\qquad
 \|e_i\|_\infty\le1+b_{i+1}\|e_{i+1}\|_\infty
                   \le1+D\|e_{i+1}\|_\infty.
\]
Backward iteration bounds all coordinates of $\ell$ by
$1+(D+2)^D$.  Also $L_i\subset p_i^{-1}\Z^{b_i}$ and the terminal
lattice is integral.  These facts give Equation~\eqref{eq:hierarchy-alphabets}.
Each ordinary scalar alphabet has at most $2Q_{\rm raw}p_i+1$ elements;
$\log p_i=O(b_i\log b_i)$ gives its asserted logarithmic bound.
\end{proof}

For the covering construction we take
$b_1=\lceil(\log\log n)^2\rceil$ once $n$ is sufficiently large.
Proposition~\ref{prop:vertical} then has
$D=O((\log\log n)^2+1)$, and both $D$ and $R$ satisfy the polynomial
bounds required in Lemma~\ref{lem:sections}.
